\section{Tooling Infrastructure}
\label{sec:tools}

The framework combines several specialized open-source tools to cover inference, web search, quality evaluation, vector storage, and neural network deployment.

\subsection{Agno Framework for Agent Orchestration}

Agno (\url{https://github.com/agno-agi/agno}) is a lightweight Python framework for building tool-augmented LLM agents. In Workflow 4, the web research node wraps an Ollama-served Mistral instance in an Agno \texttt{Agent} and equips it with \texttt{DuckDuckGoTools} for web queries. Agno handles the tool-calling loop transparently, logging each invocation for debugging. All tool calls are synchronous and run within the LangGraph node execution context.

\subsection{DeepEval for RAG Quality Assessment}

DeepEval (\url{https://github.com/confident-ai/deepeval}) implements the LLM-as-a-Judge paradigm for automated output evaluation. After Workflow 4 generates a summary, four metrics are computed against the retrieved documents:

\begin{enumerate}
    \item \textbf{Faithfulness}: $\frac{\text{Grounded Claims}}{\text{Total Claims}}$
    \item \textbf{Answer Relevancy}: $\frac{\text{Relevant Statements}}{\text{Total Statements}}$
    \item \textbf{Contextual Relevancy}: $\frac{\text{Useful Retrieved Docs}}{\text{Total Retrieved Docs}}$
    \item \textbf{Hallucination}: $1 - \frac{\text{Contradicted/Unsupported Claims}}{\text{Total Claims}}$
\end{enumerate}

All four metrics are evaluated by a local \texttt{OllamaModel} backed by \texttt{deepseek-r1:latest}, keeping the evaluation pipeline fully offline. The threshold for all metrics is set to 0.55. Evaluation runs synchronously in a separate thread via \texttt{asyncio.to\_thread} to avoid blocking the main event loop.

\subsection{ChromaDB for Vector Search}

ChromaDB (\url{https://www.trychroma.com/}) stores embedded documentation for retrieval-augmented generation in Workflow 5. Two separate vectorstores are maintained:

\begin{itemize}
    \item \textbf{Architecture Best Practices}: transfer learning guides and optimization techniques
    \item \textbf{STM32 Documentation}: hardware constraints and peripheral configurations
\end{itemize}

Embeddings are generated with \texttt{all-MiniLM-L6-v2} (384-dimensional vectors).

\subsection{Ollama: Local Development and Benchmarking}

During the initial development phase, \textbf{Ollama} (\url{https://ollama.com}) served as the primary inference backend for evaluating candidate models locally. Three practical advantages drove this choice: it kept all processing on-device (no sensitive source code or proprietary datasets left the machine), it eliminated network latency from the inference path, and it allowed unrestricted daily experimentation without per-token costs.

The benchmarking results from this phase, summarized in Section~\ref{sec:llm_selection}, informed the final model selection for the production Triton deployment.

\subsection{NVIDIA Triton Inference Server}

In the production architecture, \textbf{NVIDIA Triton Inference Server} is the primary inference backend. Centralizing model execution resolves a practical constraint encountered during development: local per-container inference engines lock the GPU for a single user, preventing concurrent requests from other workflows. Triton exposes OpenAI-compatible HTTP endpoints and manages dynamic model loading, allowing different LLMs to be swapped in and out of VRAM on demand.

Memory efficiency is addressed by \textbf{vLLM}, integrated as the serving engine behind Triton. vLLM implements \textbf{PagedAttention}, which manages the transformer KV cache as small non-contiguous physical pages rather than large pre-allocated contiguous blocks. This reduces memory fragmentation substantially and allows the system to run models with larger effective context windows within the 16~GB VRAM constraint of the development server.

\subsection{STEdgeAI CLI}

STMicroelectronics' official toolchain for neural network deployment to X-CUBE-AI:

\begin{verbatim}
stedgeai analyze --model <model.h5> --target <stm32h743>
stedgeai validate --model <model.h5> --target <stm32h743>
stedgeai generate --model <model.h5> --compression <high>
\end{verbatim}

The toolchain converts Keras or TFLite models into optimized C code compatible with the X-CUBE-AI runtime library (v9.x). The three commands are called in sequence by Workflow 2: \texttt{analyze} profiles resource usage, \texttt{validate} confirms operator support on the target MCU, and \texttt{generate} produces the final inference artifacts.

\subsection{NNI: Neural Network Intelligence}

Microsoft's NNI framework (\url{https://nni.readthedocs.io}) handles hyperparameter search in Workflow 5. Rather than using NNI's static YAML configuration files, the framework generates experiment code programmatically at runtime. This lets the search space be constructed dynamically from the actual Keras model architecture under optimization, with custom trial logic for memory footprint control and adaptive layer replacement embedded directly in the generated scripts. Best-trial extraction and retraining are also triggered programmatically at the end of each experiment.

Supported tuners include TPE (Tree-structured Parzen Estimator), Random Search, Grid Search, and SMAC.
